\section{Fixed quotients and the convention dictionary}

On $T^3$, $x,y,z$ have period $2\pi$ and $\sigma=(dx,dy,dz)$. On Nil use $\sigma=(dx,dy,dz-xdy)$ with the following lattice generators.

\stepcounter{equation}\begin{equation}\tag{QA-01}\label{eq:QA-01}
(x,y,z)\sim(x+2\pi,y,z+2\pi y)\sim(x,y+2\pi,z)\sim(x,y,z+4\pi^2).
\end{equation}

Each generator preserves the coframe, and the fundamental volume is $(2\pi)^4$. The central Fourier factor is $e^{imz/(2\pi)}$, $m\in\mathbb Z$. The dual fields $X=\partial_x$, $Y=\partial_y+x\partial_z$, $Z=\partial_z$ give $-(X^2+Y^2+Z^2)$. The $m=0$ sector is the base torus; $m\ne0$ gives an oscillator with frequency $|m|/(2\pi)$ plus the central constant $m^2/(4\pi^2)$. Twisted periodicity gives $|m|$ independent states for each signed $m$.

For Sol take $\sigma=(dx,e^xdy,e^{-x}dz)$, $A=\left(\begin{smallmatrix}2&1\\1&1\end{smallmatrix}\right)$ and $\lambda=(3+\sqrt5)/2$. Choose a unit-area fiber lattice in eigenvector coordinates of $A$ and identify $(x,y,z)$ with $(x+\log\lambda,\lambda^{-1}y,\lambda z)$. Then $n_{23}=n_{32}=1$ and the reference volume is $\log\lambda$. Fiber Fourier covectors are grouped into dual-lattice orbits under $A$.

For simply connected $S^3=SU(2)$ use the unit round coframe satisfying $d\sigma^i=2\sigma^j\wedge\sigma^k$ cyclically. The metric $q^2\sum_i(\sigma^i)^2$ has sectional curvature $q^{-2}$, scalar curvature $6q^{-2}$ and volume $2\pi^2q^3$. A structure matrix defined with a negative structure equation changes sign when converted to the present convention. Axis changes must act consistently, including orientation, on metric, connection and torsion.

\stepcounter{equation}\begin{equation}\tag{QA-02}\label{eq:QA-02}
C^i{}_{jk}=\frac{a}{r}\epsilon_{ijk}\quad\Longrightarrow\quad q=\frac{2r}{|a|},\quad R=\frac{3a^2}{2r^2},\quad {\rm Vol}=\frac{16\pi^2r^3}{|a|^3}.
\end{equation}

Equation (QA-02) converts an orthonormal-frame structure magnitude $a/r$ to physical radius. Thus $a=1$ gives $q=2r$, $a=2$ gives $q=r$, and $a=4$ gives $q=r/2$. Curvature must not be converted while volume or wave numbers are left unchanged. The present calculation uses geometric torsion $T$ directly. The supplementary comparison guide gives implementation conversions.

\stepcounter{equation}\begin{equation}\tag{QA-03}\label{eq:QA-03}
\mathcal N=\frac14[abcd]T^e{}_{ab}T_{ecd}-\frac12[abcd]R_{abcd},\qquad[0123]=+1.
\end{equation}

The alternating symbol here is not the epsilon tensor with indices raised by the Lorentz metric. Both terms of complete NY use the same connection and agree with the exterior derivative in (A-01).


\subsection{Perturbation normalizations and symbols}

In the inhomogeneous calculations a common positive action factor is suppressed and integration uses fixed Haar measure. Symmetric metric components are ordered $(11,22,33,12,13,23)$, with conjugate momentum $\delta\pi^{ii}$ on the diagonal and $2\delta\pi^{ij}$ off the diagonal. Thus $J$ is the standard canonical matrix. A dagger includes complex conjugation and the formal spatial adjoint.

\stepcounter{equation}\begin{equation}\tag{QA-04}\label{eq:QA-04}
z=(h,p),\qquad J=\begin{pmatrix}0&I\\-I&0\end{pmatrix},\qquad
\Omega=-J,\qquad [X_i,X_j]=C^k{}_{ij}X_k.
\end{equation}

On Nil, $m$ is the central Fourier integer, $r$ a lattice copy, $n$ the oscillator index and $L$ the tensor grade. On $S^3$, $j,m,n$ are Wigner labels, independent of the background Taub parameter fixed to zero. On Sol, $x$ is the base coordinate and $\xi$ the unfolded real coordinate. The Nil period is also denoted $P_0=2\pi$ to distinguish it from the Gowdy field $P$; the LRS background value of $P$ in (NG-02) is $\bar P$ and $\zeta=e^{2\bar P}$. The Sol period is $L_{\rm Sol}=\log\lambda$, distinct from the Lagrangians $L_0$ and $L^{[2]}$ and from the Nil grade $L$.
